%% kapa: exam=1 round=27 year=2016 period=2 form=A subject=accounting subject_ko=회계학 seq=37 number=77
%% source: Q-Net 감정평가사 기출문제 (cst00309) artlSeq=5169402 file=extracted/1008/5169402/제27회 감정평가사 1차 문제/제27회 감정평가사 2교시문제(A형).pdf
%% conversion: kordoc 4.24.51 -> tools/gen1.py (auto); review=proofread-sonnet; flags=none
\begin{question}{77}
(주)감평은 당기부터 단일의 공정을 거쳐 주산물 A, B, C와 부산물 X를 생산하고 있고 당기발생 결합원가는 ￦9,900이다. 결합원가의 배부는 순실현가치법을 사용하며, 부산물의 평가는 생산기준법(순실현가치법)을 적용한다. 주산물 C의 기말재고자산은?

{\small\centering
\begin{tabular}{|c|c|c|c|c|}\hline
구 분 & \makecell{최종\\생산량(개)} & \makecell{최종\\판매량(개)} & \makecell{최종 단위당\\판매가격(원)} & \makecell{추가\\가공원가(원)} \\ \hline
A & 9 & 8 & 100 & 0 \\ \hline
B & 27 & 10 & 150 & 450 \\ \hline
C & 50 & 20 & 35 & 250 \\ \hline
X & 40 & 1 & 10 & 0 \\ \hline
\end{tabular}\par}
\begin{choices}
  \item ￦800
  \item ￦1,300
  \item ￦1,575
  \item ￦1,975
  \item ￦2,375
\end{choices}
\end{question}
\answer{3}
